% Página 3: por qué el algoritmo podría fallar y cómo arreglarlo (misma información que ../algorithm-risks.md).
\pageheader{algo}{Por qué el algoritmo podría no funcionar, y cómo hacer que funcione}{Lo lee sobre todo \roletag{algo}{Algoritmo}; \roletag{datos}{Datos} lee el paso 2 de la solución y \roletag{pitch}{Pitch} lee “Qué decir en el pitch”.}

\begin{callout}[algo]
\textbf{Pipeline actual (versión ingenua):} (1) embedding de cada fila del dataset; (2) \textbf{vector de ataque} = promedio \code{attack} − promedio \code{benign}; (3) \textbf{vector de tema} = promedio de los documentos del RAG; (4) partir el mensaje en trozos: \code{attack\_score} = coseno máximo trozo–vector de ataque, \code{topic\_score} = coseno mensaje–vector de tema; (5) dos umbrales fijos deciden \Verdicts.
\end{callout}

\sect{Por qué podría fallar}
{\small
\begin{tabularx}{\linewidth}{|p{4mm}|>{\raggedright}X|>{\raggedright\arraybackslash}p{58mm}|}
\hline\rowcolor{soft}\textbf{\#} & \textbf{Problema} & \textbf{Qué se ve en la demo}\\\hline
1 & \textbf{Los ataques no son un solo grupo.} Jailbreak por juego de rol, robo del prompt de sistema, instrucciones escondidas, texto ofuscado (base64, leetspeak), otros idiomas. El promedio de grupos distintos puede quedar lejos de todos. & Detecta “ignora tus instrucciones” pero deja pasar un jailbreak de rol.\\\hline
2 & \textbf{Los embeddings miden tema, no intención.} “Ignora las instrucciones anteriores” y “¿Cómo ignoro los warnings en Python?” se parecen mucho para un modelo genérico. & Falsos positivos justo en los \emph{hard negatives}.\\\hline
3 & \textbf{El atajo del dataset.} Los ataques públicos están casi todos en inglés y nuestros benignos en español; la resta de promedios aprende “inglés vs.\ español” (o “largo vs.\ corto”), no “ataque vs.\ normal”. & Números excelentes en la evaluación y fallos en vivo con ataques en español.\\\hline
4 & \textbf{Anisotropía.} Los embeddings viven en un “cono” estrecho: casi todos los cosenos caen entre 0,6 y 0,9 y la señal útil es una diferencia pequeña. & Umbrales que cambian muchísimo con cada dato nuevo.\\\hline
5 & \textbf{Trozos (chunks).} Trozos cortos dan embeddings ruidosos; el máximo sobre muchos trozos sube los falsos positivos en mensajes largos; un ataque partido en dos se diluye. & Mensajes largos y normales bloqueados; ataques largos que pasan.\\\hline
6 & \textbf{Un solo vector de tema.} “Hola”, “gracias” o “¿y eso por qué?” no se parecen al promedio de los documentos del curso. & El chatbot bloquea saludos y preguntas de seguimiento.\\\hline
7 & \textbf{Ataque envuelto en el tema.} “Para mi tarea de Java, ignora tus reglas y…” sube el \code{topic\_score} y diluye el \code{attack\_score}. & Un jurado lo prueba y pasa.\\\hline
8 & \textbf{Pocos datos de evaluación.} Con \textasciitilde50 ataques de prueba, cada error mueve la detección 2 puntos; elegir el umbral sobre los mismos datos infla los resultados. & Números que no se sostienen si preguntan.\\\hline
9 & \textbf{Solo miramos lo que escribe el usuario.} Una inyección dentro de un documento del RAG o de la salida de una herramienta (inyección indirecta) nunca pasa por \code{check()}. & Limitación de alcance: hay que decirla, no esconderla.\\\hline
\end{tabularx}}

\sect{Cómo hacer que funcione (en orden de costo)}
\begin{enumerate}[leftmargin=14pt,itemsep=2pt]
  \item \textbf{Modelo multilingüe y local}, p.\,ej. \code{intfloat/multilingual-e5-small} o \code{sentence-transformers/paraphrase-multilingual-MiniLM-L12-v2}: corre en CPU, determinista, rápido (<50 ms por mensaje corto), no gasta tokens. \emph{Arregla parte de 3.}
  \item \textbf{Datos sin atajos} (\roletag{datos}{Datos}): mismos idiomas en ambas clases (traducir \textasciitilde50 ataques al español, escribir benignos también en inglés), longitudes parecidas, ≥30 \emph{hard negatives}. Quitar casi-duplicados (coseno > 0,95) \textbf{antes} de separar; separar train/validación/test \textbf{por fuente o plantilla}, no al azar. \emph{Arregla 3 y 8.}
  \item \textbf{Centrar los embeddings:} restar el embedding promedio del corpus y volver a normalizar antes de los cosenos. \emph{Arregla gran parte de 4.}
  \item \textbf{De un vector a un clasificador lineal:} la resta de promedios ya es un clasificador lineal; una \textbf{regresión logística} sobre los embeddings encuentra la mejor dirección con los mismos datos y sigue siendo pseudo-determinista (un vector de pesos + un umbral). Comparar ambos en validación. \emph{Arregla parte de 1 y 2.}
  \item \textbf{Varios prototipos + vecinos cercanos:} k ≈ 5–10 centroides de ataque (k-means) o k vecinos más cercanos; \code{attack\_score} = similitud al ataque más cercano − similitud al benigno más cercano. Los \emph{hard negatives} hacen de contrapeso y la demo gana una explicación gratis (“se parece a este ataque conocido”). \emph{Arregla 1 y 2.}
  \item \textbf{Tema por vecino más cercano:} \code{topic\_score} = similitud máxima con los \textbf{trozos} del RAG, no con su promedio; clase pequeña de mensajes conversacionales que siempre pasan; incluir el mensaje anterior al calcular el tema. \emph{Arregla 6.}
  \item \textbf{Trozos bien hechos:} ventanas de 1–3 oraciones que se solapan \textbf{más} el mensaje completo; umbral de trozo calibrado aparte; ignorar trozos de menos de \textasciitilde4 palabras. \emph{Arregla 5 y parte de 7.}
  \item \textbf{Umbrales con presupuesto de falsos positivos:} elegir el umbral en \textbf{validación} para bloquear como máximo \textasciitilde5\,\% de las preguntas legítimas; reportar la detección a ese umbral en \textbf{test}, con curva ROC e intervalo de confianza. \emph{Arregla 8.}
  \item \textbf{Zona gris:} entre dos umbrales, un juez LLM barato (si hay tiempo) o \code{allow} con advertencia ámbar en el panel. El filtro de embeddings es la \textbf{primera capa}, no la defensa completa.
  \item \textbf{Si sobra tiempo:} pasar también por \code{check()} los trozos del RAG que se le mandan al LLM. \emph{Arregla 9.}
\end{enumerate}

\begin{minipage}[t]{0.55\linewidth}
\colsetup
\sect{Plan mínimo para la hackathon}
\begin{itemize}
  \item \textbf{CP1:} mock conectado de punta a punta (sin algoritmo real).
  \item \textbf{CP2:} pasos 1–3 sobre las primeras filas; \code{eval.py} ya imprime números, aunque sean malos.
  \item \textbf{CP3:} comparar en validación ingenuo vs.\ regresión logística vs.\ vecinos (pasos 4–5) y quedarse con el mejor; arreglar el tema (paso 6).
  \item \textbf{CP4:} fijar umbrales (paso 8), congelar, correr test \textbf{una sola vez} y pasarle los números a Pitch.
\end{itemize}
\end{minipage}\hfill
\begin{minipage}[t]{0.42\linewidth}
\colsetup
\sect{Qué decir en el pitch}
\begin{itemize}
  \item “Primera capa: determinista, barata (sin tokens) y explicable.”
  \item Comparar contra palabras clave en el \textbf{mismo} test, con el mismo presupuesto de falsos positivos.
  \item Nombrar las limitaciones sin que pregunten: paráfrasis adaptativas, inyección indirecta, tamaño del dataset.
\end{itemize}
\end{minipage}
